\section{Agentic Retrieval and Optimization}
\label{sec:optimization}

The core innovation of VulnRAG lies in its ability to actively optimize the retrieval process through agentic self-correction and structural structural discovery.

\subsection{Zero-Shot Self-Correction Loop}
To ensure retrieval reliability, we implement a training-free feedback loop (Fig.~\ref{fig:self_correction}). The `self_correction.py` module acts as an autonomous quality controller:
\begin{itemize}
    \item \textbf{Evaluation Phase}: After an initial retrieval pass, the \textit{Evaluator Agent} analyzes the result set across three dimensions: result count, mean semantic confidence score, and metadata diversity.
    \item \textbf{Reformulation Phase}: If the retrieval quality is deemed insufficient (e.g., zero results or low confidence), the `query_reformulator.py` module generates an expanded query. This process includes automated \textbf{CPE Extraction} and security synonym mapping (e.g., mapping "RCE" to "Remote Code Execution").
\end{itemize}

\subsection{Structural Discovery via CVE Relationship Graphs}
Beyond text-based similarity, VulnRAG leverages structural relationships to discover attack chains. The `cve_graph.py` service constructs a Directed Acyclic Graph (DAG) where nodes represent CVEs and edges represent explicit functional links (Fig.~\ref{fig:graph}).
\begin{itemize}
    \item \textbf{Edge Extraction}: Edges are derived from URL references in NVD advisories, shared CWE (Common Weakness Enumeration) identifiers, and product-vendor clusters.
    \item \textbf{Graph Traversal}: When a high-priority CVE is identified, the system performs a traversal with configurable depth (1–3 hops) to surface "hidden" related vulnerabilities that may be semantically distinct but operationally linked.
\end{itemize}

\section{Multi-Factor Semantic Reranking (MFSR)}
\label{sec:reranking}
To mitigate the "Lost in the Middle" problem \cite{lostmiddle}, we implement a Multi-Factor Semantic Reranking (MFSR) algorithm in `cve_reranker.py`. 

\subsection{Operational Relevance Formalization}
The final rank of a retrieved CVE $c$ for query $q$ is determined by the Operational Relevance Score $S_{OR}$:
\begin{equation}
S_{OR}(q, c) = w_{sem} S_{sem}(q, c) + w_{cvss} V_{cvss}(c) + w_{exp} E_{exp}(c) + w_{rec} R_{rec}(c)
\label{eq:mfsr}
\end{equation}
where $S_{sem}$ is the semantic affinity from a Cross-Encoder (\texttt{BGE-reranker-large}), and $V_{cvss}$, $E_{exp}$, and $R_{rec}$ are normalized domain scores for severity, exploitability, and recency, respectively.

\section{Complexity-Aware Multi-Tier Routing}
\label{sec:routing}
To ensure economic feasibility, VulnRAG implements a tiered execution model. The `SmartCVESearchService` dynamically routes sub-tasks to different LLM tiers based on complexity indicators:
\begin{itemize}
    \item \textbf{Flash Tier}: Routine NVD metadata lookups.
    \item \textbf{Pro Tier}: Complex query reformulation and entity extraction.
    \item \textbf{Reasoning Tier}: Deep intel synthesis and attack path modeling.
\end{itemize}
Escalation occurs only when the Evaluator detects a high-complexity reasoning failure, ensuring that flagship models are used judiciously.
